%%%%%%%%%%%%%%%%%%%%%%%%%%%%%%%%%%%%%%%%%%%%%%%%%%%%%%%%%%%%%%%%%%%%%%
% Lớp 10 | Xác suất | Không gian mẫu và biến cố
% Dạng 0D0?1-2: Mô tả không gian mẫu, biến cố 
% Số câu: 67  (sinh tự động bởi tach_cau.py; từ nay thầy sửa trực tiếp file này)
%%%%%%%%%%%%%%%%%%%%%%%%%%%%%%%%%%%%%%%%%%%%%%%%%%%%%%%%%%%%%%%%%%%%%%

%%%%%%%%%%%%%%%% TRẮC NGHIỆM 4 PHƯƠNG ÁN %%%%%%%%%%%%%%%%

% nguồn: MaHoaNH_TracNghiem_Lop10
\begin{ex}%[0D0N1-2]%[ma:c9feda72]
	Gieo một con súc sắc cân đối và đồng chất hai lần. Hãy phát biểu biến cố $ A=\{(6, 1); (6, 2);  (6, 3); (6,4); (6, 5); (6, 6)\} $ dưới dạng mệnh đề.
	 \choice
		{$  A\colon$\lq\lq Tổng số chấm xuất hiện lớn hơn $ 6 $\rq\rq}
		{$ A\colon$\lq\lq Mặt $ 6 $ chấm  xuất hiện\rq\rq}
		{\True $ A\colon$\lq\lq Lần đầu xuất  hiện mặt $ 6 $ chấm\rq\rq}
		{$ A\colon$\lq\lq Tổng số chấm không nhỏ hơn $ 7  $\rq\rq}
	 \loigiai{
		Biến cố $ A $ được phát biểu dưới dạng mệnh đề là \lq\lq Lần  đầu xuất hiện mặt $ 6 $ chấm\rq\rq.
	 }
\end{ex}

% nguồn: MaHoaNH_TracNghiem_Lop10
\begin{ex}%[0D0N1-2]%[ma:3bd5c2e7]
	Gieo hai đồng tiền một lần. Kí hiệu $S$, $N$ lần lượt để chỉ đồng tiền xuất hiện mặt sấp,  đồng tiền xuất hiện mặt ngửa. Mô tả không gian mẫu nào dưới đây là đúng?
	\choice 
		{$\Omega=\{S;  N\}$}
		{$\Omega=\{NN; SS\}$} 
		{$\Omega=\{SN; NS\}$} 
		{\True  $\Omega=\{SN; NS; SS; NN\}$}
	\loigiai{ 
		Không  gian mẫu $\Omega=\{SN; NS; SS; NN\}$.
	} 
\end{ex}

% nguồn: MaHoaNH_TracNghiem_Lop10
\begin{ex}%[0D0N1-2]%[ma:982ad5d9]
	Gieo con súc sắc hai lần. Biến cố nào dưới  đây là biến cố tổng số chấm trong hai lần gieo bằng $7$?
	\choice 
		{\True $A = \left\lbrace (1;6), (6;1), (3;4), (4;3), (5;2),  (2;5)\right\rbrace $}
		 {$B = \left\lbrace (1;6), (3;4), (5;2)\right\rbrace $}
		 {$C = \left\lbrace (1;6), (6;1), (3;4), (4;3)\right\rbrace $}
		{$D = \left\lbrace (2;6), (6;2), (3;4),  (4;3), (5;2), (2;5)\right\rbrace $}
	 \loigiai{
		Biến cố tổng số chấm trong hai lần gieo súc sắc bằng $7$ là $A =  \left\lbrace (1;6), (6;1), (3;4), (4;3), (5;2), (2;5)\right\rbrace $.
	 }
\end{ex}

% nguồn: MaHoaNH_TracNghiem_Lop10
\begin{ex}%[0D0N1-2]%[ma:76a9632f]
	Gieo một con  xúc sắc cân đối và đồng chất hai lần. Biến cố $A$: \lq\lq Lần đầu xuất hiện mặt $5$ chấm\rq\rq\ là
	 \choice
		{\True  $A=\{(5;1),(5;2),(5;3),(5;4),(5;5),(5;6)\}$}
		 {$A=\{(5;1)\}$}
		 {$A=\{(1;5),(2;5),(3;5),(4;5),(5;5),(6;5)\}$}
		{$A=\{(5;1),(5;2),(5;3),(5;4),(5;5),(5;6),(1;5),(2;5),(3;5),(4;5),(5;5),(6;5)\}$} 
	\loigiai{ 
		Biến cố $A$:  \lq\lq Lần đầu xuất hiện mặt $5$ chấm\rq\rq\ là $A=\{(5;1),(5;2),(5;3),(5;4),(5;5),(5;6)\}$.
	 }
\end{ex}

% nguồn: MaHoaNH_TracNghiem_Lop10
\begin{ex}%[0D0N1-2]%[ma:bedec663]
	Gieo $2$ đồng tiền là một phép thử ngẫu nhiên có không gian mẫu là 
	\choice 
		{\True $\left\lbrace NN; NS; SN; SS\right\rbrace  $}
		{$\left\lbrace  NS; SN\right\rbrace $}
		{$\left\lbrace  N; S\right\rbrace $}
		{$\left\lbrace NNN; SSS; NNS; SSN; NSS; SNN\right\rbrace  $}
	 \loigiai{
		Gieo $2$ đồng tiền là một phép thử ngẫu nhiên có không gian mẫu là $\Omega  = \left\lbrace NN; NS; SN; SS\right\rbrace$.
	} 
\end{ex}

% nguồn: MaHoaNH_TracNghiem_Lop10
\begin{ex}%[0D0N1-2]%[ma:f25f4b05]
	Từ một hộp chứa $6$ quả cầu trắng, $7$ quả cầu đỏ, $4$ quả  cầu vàng, các quả cầu có kích thước và khối lượng giống nhau. Xét phép thử \lq \lq Lấy ngẫu nhiên đồng thời $3$ quả cầu\rq \rq. Biến cố đối của biến cố $A \colon$ \lq \lq $3$ quả cầu lấy ra có ít nhất một quà màu vàng\rq \rq \ là
	\choice 
		{\True $\overline{A} \colon$ \lq  \lq
		$3$ quả cầu lấy ra không có màu vàng\rq \rq} 
		{$\overline{A}\colon$ \lq \lq $3$ quả cầu lấy ra cùng màu trắng\rq  \rq}
		 {$\overline{A}\colon$ \lq \lq $3$ quả cầu lấy ra cùng màu đỏ\rq \rq}
		{$\overline{A}\colon$ \lq \lq $3$  quả cầu lấy ra có $3$ màu khác nhau\rq \rq}
	 \loigiai{
		Biến cố đối của biến cố $A \colon$ \lq \lq $3$ quả  cầu lấy ra có ít nhất một quà màu vàng\rq \rq \ là $\overline{A} \colon$ \lq \lq
		$3$ quả  cầu lấy ra không có màu vàng\rq \rq.
	} 
\end{ex}

% nguồn: MaHoaNH_TracNghiem_Lop10
\begin{ex}%[0D0N1-2]%[ma:bc8bab69]
	 Trong một hộp có
	$2$ viên bi trắng được đánh số từ $1$ đến  $2$;
	$3$ viên bi xanh được đánh  số từ $3$ đến $5$;
	$2$  viên bi đỏ được đánh số từ $6$ đến $7$.
	Lấy ngẫu nhiên hai viên bi, mô tả  không gian mẫu nào dưới đây là đúng?
	 \choice
		{$\Omega =  \left\{\left(m, n\right)\mid 1 \le m \le 7, 1 \le n \le 7\right\}$}
		{$\Omega = \left\{\left(m, n\right)\mid 1 \le m \le 5, 6  \le n \le 7\right\}$}
		{\True $\Omega = \left\{\left(m, n\right)\mid 1 \le m \le 7, 1  \le n \le 7, m \ne n\right\}$}
		{$\Omega = \left\{\left(m, n\right)\mid 1  \le m \le 3, 4 \le n \le 4\right\}$}
	 \loigiai{
		Mỗi viên bi đánh một số, nên  $2$ viên bi lấy ra sẽ có số khác nhau.\\
		$\Omega = \left\{\left(m, n\right)\mid 1 \le m  \le 7, 1 \le n \le 7, m \ne n\right\}$.
	} 
\end{ex}

% nguồn: MaHoaNH_TracNghiem_Lop10
\begin{ex}%[0D0H1-2]%[ma:9d35bc14]
	Gieo một đồng xu cân đối, đồng chất $3$ lần là  một phép thử ngẫu nhiên có không gian mẫu là
	\choice 
		{$\{NN,SN,NS,SS\}$} 
		{\True  $\{NNN,SSS,NNS,SSN,NSN,SNS,NSS,SNN\}$}
		 {$\{NNN,SSS,NNS,SSN,NSN,SNS\}$}
		 {$\{NNN,SSS,NNS,SSN,NSS,SNN\}$}
	\loigiai{Không gian  mẫu là $\{NNN,SSS,NNS,SSN,NSN,SNS,NSS,SNN\}$.}
\end{ex}

% nguồn: MaHoaNH_TracNghiem_Lop10
\begin{ex}%[0D0H1-2]%[ma:ed702e46]
	Gieo $3$ đồng tiền cân đối đồng chất  là một phép thử ngẫu nhiên có không gian mẫu là
	 \choice
		 {$\{NNN,SSS,NNS,SSN,NSN,SNS\}$}
		 {$\{SN,NS,SS,NN\}$}
		{\True $\{NNN,SSS,NNS,SSN,NSN,SNS,NSS,SNN\}$} 
		 {$\{NNN,SSS,NNS,SSN,NSS,SNN\}$}
	\loigiai{Khi gieo $3$ đồng tiền cân đối, đồng chất, phép thử ngẫu  nhiên này có không gian mẫu là $$\{NNN,SSS,NNS,SSN,NSN,SNS,NSS,SNN\}.$$
	} 
\end{ex}

% nguồn: MaHoaNH_TracNghiem_Lop10
\begin{ex}%[0D0H1-2]%[ma:81dca4cc]
	Gieo một đồng xu cân đối và đồng chất liên tiếp ba lần. Gọi $A$ là biến cố \lq\lq Có ít nhất hai mặt sấp xuất hiện liên tiếp\rq\rq\ và $B$ là  biến cố \lq\lq Kết quả ba lần gieo là như nhau\rq\rq. Xác định biến cố $A\cup B$.
	\choice 
		{$A\cup B=\{SSS,SSN,NSS,SNS,NNN\}$} 
		{$A\cup B=\{SSS,NNN\}$} 
		 {\True $A\cup B=\{SSS,SSN,NSS,NNN\}$}
		{$A\cup B=\Omega$} 
	\loigiai{Ta  có $A=\{SSS,SSN,NSS\}$, $B=\{SSS,NNN\}$.\\
	Vậy $A\cup  B=\{SSS,SSN,NSS,NNN\}$.}
\end{ex}

% nguồn: MaHoaNH_TracNghiem_Lop10
\begin{ex}%[0D0H1-2]%[ma:0dee273b]
	Gieo một đồng xu cân đối đồng chất $3$ lần. Gọi $A_i$ là biến cố \lq\lq mặt sấp xuất hiện lần gieo thứ $i$\rq\rq, với $i=1,2,3$. Khi đó, biến cố $\overline{A}_1\cup \overline{A}_2\cup  \overline{A}_3$ là
	\choice 
		{\lq\lq  Cả $3$ lần gieo đều được mặt sấp\rq\rq}
		{\lq\lq Mặt sấp xuất hiện không quá  một lần\rq\rq}
		{\True \lq\lq Mặt ngửa xuất hiện ít nhất một lần\rq\rq} 
		{\lq\lq  Cả $3$ lần gieo đều được mặt ngửa\rq\rq}
	\loigiai{Từ $\overline{A}_1\cup \overline{A}_2\cup \overline{A}_3$ tức là hoặc mặt ngửa xuất hiện ở lần gieo thứ nhất hoặc mặt ngửa xuất hiện ở lần gieo thứ hai hoặc mặt ngửa xuất hiện ở lần gieo  thứ ba.\\
	 Vậy mặt ngửa xuất hiện ít nhất một lần.}
\end{ex}

% nguồn: MaHoaNH_TracNghiem_Lop10
\begin{ex}%[0D0H1-2]%[ma:c5f8b114]
	Gieo $3$ đồng tiền là một phép  thử ngẫu nhiên có không gian mẫu là
	 \choice
		 {$\{NN,NS,SN,SS\}$}
		 {$\{NNN,SSS,NNS,SSN,NSN,SNS\}$}
		{\True $\{NNN,SSS,NNS,SSN,NSN,SNS,NSS,SNN\}$} 
		 {$\{NNN,SSS,NNS,SSN,NSS,SNN\}$}
	 \loigiai{
		Gieo $3$ đồng tiền là một phép thử ngẫu nhiên có không gian mẫu là $$\Omega=\{NNN,SSS,NNS,SSN,NSN,SNS,NSS,SNN\}.$$ 
	 }
\end{ex}

% nguồn: MaHoaNH_TracNghiem_Lop10
\begin{ex}%[0D0H1-2]%[ma:6238b190]
	Xét phép thử, gieo một đồng tiền xu cân đối và đồng  chất ba lần. Biến cố $A$: \lq\lq Cả ba lần gieo cùng sấp hoặc cùng ngửa\rq\rq. Biến cố $A$ biểu diễn dưới dạng tập hợp là
	\choice 
		 {\True $A=\left\{SSS;NNN \right\}$}
		 {$A=\left\{NSS;NNN \right\}$}
		 {$A=\left\{SSS;NNS \right\}$}
		{$A=\left\{SNS;NSN  \right\}$}
	\loigiai{ 
		$A=\left\{SSS;NNN \right\}$ 
	} 
\end{ex}

% nguồn: MaHoaNH_TracNghiem_Lop10
\begin{ex}%[0D0H1-2]%[ma:78774703]
	Gieo ngẫu nhiên một đồng tiền cân đối và đồng chất $5$ lần. Tính số phần tử  không gian mẫu.
	\choice 
		 {$10$}
		{\True  $32$}
		{$16$} 
		 {$64$}
	\loigiai{ 
		 Ta có $n\left(\Omega\right)=2^5=32$.
	} 
\end{ex}

% nguồn: MaHoaNH_TracNghiem_Lop10
\begin{ex}%[0D0H1-2]%[ma:6b2e4407]
	Gieo một con súc sắc cân đối và đồng chất hai lần. Hãy mô tả biến cố $A$: \lq\lq Lần  đầu tiên xuất hiện mặt năm chấm\rq\rq.
	 \choice
		 {$A=\{5\}$}
		 {$A=\{5;5\}$}
		 {$A=\{(5;1);(5;2);(5;3);(5;4);(5;6)\}$}
		 {\True $A=\{(5;1);(5;2);(5;3);(5;4);(5;5);(5;6)\}$}
	\loigiai{ 
		Biến cố  \lq\lq Lần đầu tiên xuất hiện mặt năm chấm\rq\rq\ gồm các phần tử là
		$$\{(5;1),(5;2),(5;3),(5;4),(5;5),(5;6)\}.$$ 
	} 
\end{ex}

% nguồn: MaHoaNH_TracNghiem_Lop10
\begin{ex}%[0D0H1-2]%[ma:5554ec25]
	Rút ngẫu nhiên một hộp chứa $9$ tấm thẻ được đánh số $1$;$2$;$3$;$4$;$5$;$6$;$7$;$8$;$9$. Gọi $A$ là biến cố: \lq\lq Rút được thẻ ghi số chia  hết cho 3\rq\rq. Biến cố đối của biến cố $A$ là
	 \choice
		{\True $\overline{A} =  \left\lbrace 1;2;4;5;7;8\right\rbrace $}
		{$\overline{A} = \left\lbrace 3;6;9\right\rbrace  $}
		{$\overline{A} = \left\lbrace 1;2;4;5;7;8;10\right\rbrace $} 
		{$\overline{A} = \left\lbrace 2;4;6;8\right\rbrace $} 
	 \loigiai{
		Vì $A$ là biến cố: \lq\lq Rút  được thẻ ghi số chia hết cho 3\rq\rq \\
		$\Rightarrow$ biến cố đối của $A$ là $\overline{A}\colon $ \lq\lq Rút được thẻ ghi số không chia hết  cho 3\rq\rq, do đó
		$$\overline{A} = \left\lbrace  1;2;4;5;7;8\right\rbrace.$$
	} 
\end{ex}

% nguồn: MaHoaNH_TracNghiem_Lop10
\begin{ex}%[0D0H1-2]%[ma:f601c4fb]
	Hai  xạ thủ cùng bắn vào bia. Xét các biến cố sau:\\
	$A$: \lq\lq Không ai  bắn trúng\rq\rq;\\
	$B$: \lq\lq Cả  hai đều bắn trúng\rq\rq;\\
	$C$: \lq\lq Có đúng nhất một người bắn trúng\rq\rq;\\ 
	$D$: \lq\lq Có ít nhất một người bắn  trúng\rq\rq.\\
	Chọn  phát biểu {\bf sai}.
	\choice 
		 {$A=\overline{D}$}
		 {\True $B\cup C=\Omega$}
		{$D=\overline{A}$} 
		 {$B\cap C=\varnothing$}
	\loigiai{ 
		Kí hiệu $T$; $K$ lần lượt là xạ thủ bắn trúng và bắn không trúng  bia.
		Ta  có:\\
		$\Omega = \{TT; TK; KT;  KK\}$.\\
		 $A = \{KK\}$.\\
		$B = \{TT\}$.\\ 
		$C = \{TK; KT\}$.\\ 
		$D = \{TK; KT; TT \}$ 
		Do  đó $B \cup C = \{TK; KT; TT \}$.\\
		 Vậy $B\cup C=\Omega$ là phát biểu sai.
	} 
\end{ex}


%%%%%%%%%%%%%%%% ĐÚNG/SAI %%%%%%%%%%%%%%%%

% nguồn: MaHoaNH_DungSai_Lop10
\begin{ex}%[0D0N1-2]%[ma:4dca26b9]
	Xét phép thử gieo một đồng tiền hai  lần với các biến cố: $A$: \lq\lq Kết quả hai lần gieo là như nhau\rq\rq, $B$: \lq\lq Có ít nhất một lần xuất hiện mặt sấp\rq\rq, $C$: \lq\lq Lần thứ hai xuất hiện mặt sấp\rq\rq, $D$: \lq\lq Không xuất hiện mặt ngửa\rq\rq. Xét tính đúng sai của các khẳng định sau:
	 \choiceTF
		 {\True $n(A)=2$}
		 {$n(B)=2$}
		 {\True $n(C)=2$}
		 {$n(D)=2$}
	 \loigiai{
		\begin{itemchoice}
			\itemch  Ta có $A=\{SS,NN\}$, suy ra $n(A)=2$.
			 \itemch Ta có $B=\{SN,NS,SS\}$, suy ra $n(B)=3$.
			\itemch Ta có $C=\{NS,SS\}$,  suy ra $n(C)=2$.
			\itemch Ta có $D=\{SS\} $, suy  ra $n(D)=1$.
		\end{itemchoice}
	 }
\end{ex}

% nguồn: MaHoaNH_DungSai_Lop10
\begin{ex}%[0D0N1-2]%[ma:05bb2b39]
	Gieo ngẫu nhiên một  con súc sắc cân đối và đồng chất. Biết các biến cố sau
	\begin{itemize}
		\item Biến cố $A=\{1;3;5\}$. 
		\item  Biến cố $B=\{3;4;5;6\}$.
		\item Biến cố $C=\{3;6\}$. 
		\item Biến  cố $D=\{1;2;3;4\}$.
	\end{itemize}
	Hãy chọn  các khẳng định đúng?
	\choiceTF 
		 {\True $A$: "mặt lẻ xuất hiện"}
		{$B$: "xuất hiện mặt  có số chấm lớn hơn $3$"}
		{\True $C$: "xuất hiện mặt có số chấm chia hết cho  $3$"}
		{$D$: "xuất hiện  mặt có số chấm chia không quá $5$"}
	 \loigiai{
		\begin{itemchoice}
			\itemch $A=\{1;3;5\}$: biến cố mặt lẻ xuất hiện là khẳng  định đúng.
			 \itemch $B=\{3;4;5;6\}$: biến cố xuất hiện mặt có số chấm lớn hơn $3$ là sai.
			\itemch $C=\{3;6\}$: biến cố xuất  hiện mặt có số chấm chia hết cho $3$ là đúng.
			\itemch $D=\{1;2;3;4\}$: biến cố xuất hiện mặt có số chấm chia không quá $5$ là sai  vì thiếu mặt $6$ chấm.
		\end{itemchoice}
	 }
\end{ex}

% nguồn: MaHoaNH_DungSai_Lop10
\begin{ex}%[0D0N1-2]%[ma:fa8f453b]
	Phép thử gieo ngẫu nhiên một con súc sắc cân đối, đồng chất, sau đó  gieo tiếp một đồng xu có hai mặt sấp (ký hiệu S), ngữa (ký hiệu N). Biết các biến cố sau
	\begin{itemize}
		\item Biến cố  $A=\{S1;S2;S3,S4;S5;S6\}$.
		\item  Biến cố $B=\{1S;2S;3S,4S;5S;6S\}$.
		 \item Biến cố $C=\{1S;3S;3S,1N;3N;5N\}$.
		\item  Biến cố $D=\{1;2;3;4;5;6;S;N\}$.
	\end{itemize}
	Hãy chọn các khẳng định đúng? 
	 \choiceTF
		 {\True $B$: ``Mặt sấp luôn xuất hiện''}
		{$A$: ``Mặt sấp luôn xuất hiệ''} 
		 {\True $C$: ``súc sắc luôn xuất hiện mặt lẻ''}
		{$D=\Omega$: là  không gian mẫu của phép thử}
	\loigiai{ 
		\begin{itemize}
			 \item Biến cố $B=\{1S;2S;3S,4S;5S;6S\}$ là biến cố ``Mặt sấp luôn xuất hiện''.
			\item Biến cố $A=\{S1;S2;S3,S4;S5;S6\}$ không  có vì giao súc sắc rồi gieo tiếp đồng xu.
			\item Biến cố $C=\{1S;3S;3S,1N;3N;5N\}$ là  biến cố ``súc sắc luôn xuất hiện mặt lẻ''.
			\item Biến cố $D=\{1;2;3;4;5;6;S;N\}$ không phải là không gian mẫu  vì $\Omega=\{1S;2S;3S,4S;5S;6S;1N;2N;3N,4N;5N;6N\}$.
		\end{itemize}
	} 
\end{ex}

% nguồn: MaHoaNH_DungSai_Lop10
\begin{ex}%[0D0N1-2]%[ma:076ee0c3]
	Gọi $S$ là tập hợp các số tự nhiên có ba chữ số. Chọn ngẫu nhiên một số từ tập $S$. Xét tính đúng sai của các  khẳng định sau:
	 \choiceTF
		{$n(\Omega)=1\,000$} 
		{\True Gọi $A$ là biến cố: \lq\lq Chọn được số tự nhiên có các chữ  số đôi một khác nhau\rq\rq, khi đó $n(A)=648$}
		{\True Gọi $B$ là biến cố: \lq\lq Chọn được số tự nhiên chia hết cho $5$\rq\rq,  khi đó $n(B)=180$}
		{Gọi $C$ là biến cố: \lq\lq Chọn được số tự nhiên chẵn\rq\rq, khi đó  $n(C)=500$}
	 \loigiai{
		\begin{itemchoice}
			 \itemch Xét số tự nhiên có ba chữ số dạng $\overline{abc}$.\\
			Số cách chọn $a$ ($a$ khác $0$) và $b, c$ lần lượt là $9$, $10$, $10$ nên số  các số tự nhiên gồm ba chữ số là $9 \cdot 10 \cdot 10 = 900$.\\
			Phép thử đang xét là hoạt động chọn ngẫu nhiên một số từ $S$ nên số kết quả thuận lợi  không gian mẫu là $n(\Omega) = 900$.
			\itemch Xét số tự nhiên có ba chữ số dạng  $\overline{abc}$.
			\begin{itemize}
				\item Chọn $a$ ($a \ne 0$):  có $9$ cách.
				\item Chọn  $b$ ($b \ne a$): có $9$ cách.
				\item Chọn $c$ ($c \ne a$, $c \ne b$): có $8$ cách. 
			\end{itemize}
			Vậy số các số tự nhiên có ba chữ số  đôi một khác nhau là $9 \cdot 9 \cdot 8 = 648$.\\
			Vì  vậy $n(A) = 648$.
			\itemch Xét số tự nhiên có ba chữ số  dạng $\overline{abc}$.\\
			Số  này chia hết cho $5$ nên $c \in \{0;5\}$: có $2$ cách chọn $c$.\\
			Số cách chọn $a$  ($a$ khác $0$), $b$ lần lượt là $9$, $10$.\\
			Vậy số các số tự nhiên thỏa mãn là $2  \cdot 9 \cdot 10 = 180$.\\
			Vì  vậy $n(B) = 180$.
			\itemch Xét số tự  nhiên có ba chữ số dạng $\overline{abc}$.\\
			Số này là số chẵn nên $c$ có $5$ cách chọn.\\ 
			$a$ có $9$ cách chọn, $b$ có  $10$ cách chọn.\\
			Vậy số các số tự nhiên thỏa mãn là $9  \cdot 10 \cdot 5 = 450$.
		\end{itemchoice}
	} 
\end{ex}

% nguồn: MaHoaNH_DungSai_Lop10
\begin{ex}%[0D0H1-2]%[ma:5b4863bc]
	Gieo hai con xúc xắc. Xét tính  đúng sai của các khẳng định sau:
	 \choiceTF
		{\True Số kết  quả thuận lợi cho biến cố \lq\lq Số chấm trên hai con hơn kém nhau $2$\rq\rq\, bằng $8$}
		{Số kết quả thuận  lợi cho biến cố \lq\lq Tích số chấm chia hết cho $5$\rq\rq\, bằng $12$}
		{\True Số kết quả thuận lợi cho biến cố \lq\lq Tích số chấm là  số lẻ\rq\rq\, bằng $9$}
		{Số kết quả thuận lợi cho biến cố \lq\lq Tổng số chấm là số chẵn\rq\rq\,  bằng $15$}
	 \loigiai{
		\begin{itemchoice}
			\itemch Các cặp: $(1,3)$, $(2,4)$, $(3,5)$, $(4,6)$ và đối xứng của  chúng $\Rightarrow 8$ kết quả.
			\itemch Các cặp có tích chia hết cho $5$ phải có một số  là $5$: có $11$ cặp. (Không có cặp $(5,5)$ lặp).
			\itemch Cả hai số đều lẻ: $1,3,5$ $\Rightarrow 3 \cdot 3 =  9$.
			\itemch Tổng chẵn khi cả hai cùng chẵn ($2,4,6$) hoặc cùng lẻ ($1,3,5$): $3 \cdot 3 + 3 \cdot  3 = 18$.
		\end{itemchoice}
	 }
\end{ex}

% nguồn: MaHoaNH_DungSai_Lop10
\begin{ex}%[0D0H1-2]%[ma:81f3e695]
	Gọi $A$ là tập hợp các số  tự nhiên có $2$ chữ số nhỏ hơn $20$. Lấy ra $1$ số tự nhiên bất kỳ trong $A$. Xét tính đúng sai của các khẳng định sau:
	\choiceTF 
		 {\True $n(\Omega)=10$}
		{\True Gọi  $B$ là biến cố: \lq\lq Lấy được một số tự nhiên lẻ\rq\rq. Khi đó $n(B)=5$}
		{Gọi $C$ là biến cố: \lq\lq Lấy được một số tự nhiên chia hết cho $3$\rq\rq. Khi đó $n(C)=2$} 
		{Gọi $D$ là biến cố: \lq\lq Lấy được một số  nguyên tố\rq\rq. Khi đó $n(D)=3$}
	\loigiai{ 
		\begin{itemchoice}
			\itemch Ta có không gian mẫu: $\Omega = \{10,11,12,13,14,15,16,17,18,19\} \Rightarrow n(\Omega) =  10$.
			 \itemch Ta có $B = \{11,13,15,17,19\} \Rightarrow n(B) = 5$.
			\itemch Ta có  $C = \{12,15,18\} \Rightarrow n(C) = 3$.
			 \itemch Ta có $D = \{11,13,17,19\} \Rightarrow n(D) = 4$.
		\end{itemchoice}
	} 
\end{ex}

% nguồn: MaHoaNH_DungSai_Lop10
\begin{ex}%[0D0H1-2]%[ma:eb6f25e5]
	Trong hộp có $3$ bi xanh, $4$ bi đỏ và $5$ bi vàng có kích thước và khối lượng như nhau. Lấy ngẫu nhiên từ trong hộp $4$ viên bi. Xét  tính đúng sai của các khẳng định sau:
	\choiceTF 
		{\True Số phần tử  của không gian mẫu bằng $495$}
		{\True Số các kết quả thuận  lợi cho biến cố \lq\lq Trong $4$ viên bi được chọn có ít nhất $1$ bi xanh\rq\rq\, bằng $369$}
		{Số  các kết quả thuận lợi cho biến cố \lq\lq Trong $4$ viên bi được chọn có đúng $1$ viên bi đỏ\rq\rq\, bằng $220$}
		{Số các kết quả thuận lợi cho biến  cố \lq\lq Trong $4$ viên bi được chọn có ít nhất $2$ bi đỏ\rq\rq\, bằng $199$}
	 \loigiai{
		\begin{itemchoice}
			\itemch Ta có không gian mẫu là  kết quả của tổ hợp chập $4$ của $12$ phần tử, do đó số phần tử của không gian mẫu là $n(\Omega) = \mathrm{C}_{12}^4 = 495$.
			\itemch Số kết quả thuận lợi cho biến cố \lq\lq Có ít nhất $1$ bi xanh\rq\rq\, là số phần tử của $n(\Omega) - \mathrm{C}_9^4$, trong đó $\mathrm{C}_9^4$ là số cách chọn $4$ viên bi không có bi xanh từ $9$ viên còn lại (bi đỏ và bi vàng).\\ Ta có $n(\Omega) -  \mathrm{C}_9^4 = 495 - 126 = 369$.
			\itemch Số kết quả thuận lợi cho biến cố \lq\lq Có đúng  $1$ viên bi đỏ\rq\rq\, là số cách chọn $1$ bi đỏ từ $4$ viên bi đỏ và $3$ viên bi xanh, đồng thời chọn $3$ bi vàng từ $5$ viên bi vàng.\\ Ta có $\mathrm{C}_4^1 \cdot \mathrm{C}_5^3 = 4 \cdot 10 = 40$ chứ không phải $220$.
			\itemch Số kết quả thuận lợi cho biến cố \lq\lq Có ít nhất $2$ bi đỏ\rq\rq\, là số phần tử của $n(\Omega) - \mathrm{C}_8^4 - \mathrm{C}_4^1 \cdot \mathrm{C}_8^3$, trong đó $\mathrm{C}_8^4$ là số cách chọn  $4$ viên bi không có bi đỏ (chỉ chọn bi xanh và bi vàng), còn $\mathrm{C}_4^1 \cdot \mathrm{C}_8^3$ là số cách chọn $1$ bi đỏ và $3$ bi không đỏ.\\ Ta có $495 - 70 - 160 = 265$ chứ không phải $199$.
		\end{itemchoice}
	} 
\end{ex}

% nguồn: MaHoaNH_DungSai_Lop10
\begin{ex}%[0D0H1-2]%[ma:2b10ed99]
	Xét phép thử là gieo một con xúc xắc một lần. Xét tính đúng sai của các khẳng  định sau:
	 \choiceTF
		 {\True $n(\Omega)=6$}
		{\True Số kết quả thuận lợi của biến cố: \lq\lq Thu được mặt có số chấm chia hết cho $2$\rq\rq\, bằng  $3$}
		 {Số kết quả thuận lợi của biến cố: \lq\lq Thu được mặt có số chấm nhỏ hơn $5$\rq\rq\, bằng $3$}
		{Số kết quả thuận lợi của biến cố: \lq\lq Thu được  mặt có số chấm là số lẻ\rq\rq\, bằng $4$}
	\loigiai{ 
		\begin{itemchoice}
			\itemch Ta có $\Omega = \{1,2,3,4,5,6\} \Rightarrow n(\Omega) = 6$. 
			\itemch Ta  có $A = \{2,4,6\} \Rightarrow n(A) = 3$.
			\itemch Ta có $B =  \{1,2,3,4\} \Rightarrow n(B) = 4$.
			\itemch Ta  có $C = \{1,3,5\} \Rightarrow n(C) = 3$.
		\end{itemchoice}
	 }
\end{ex}

% nguồn: MaHoaNH_DungSai_Lop10
\begin{ex}%[0D0H1-2]%[ma:50477e0e]
	Gieo $5$ lần một đồng tiền hai mặt sấp, ngửa. Xét tính đúng  sai của các khẳng định sau:
	\choiceTF 
		{\True  $n(\Omega)=32$}
		{\True Số kết quả thuận lợi của biến cố $A$: \lq\lq Lần đầu  tiên xuất hiện mặt ngửa\rq\rq\, bằng $16$}
		{Số kết quả thuận lợi của biến cố $B$: \lq\lq Mặt sấp xuất hiện ít nhất một lần\rq\rq\, bằng  $30$}
		{\True  Số kết quả thuận lợi của biến cố $C$: \lq\lq Số lần mặt sấp xuất hiện nhiều hơn mặt ngửa\rq\rq\, bằng $16$}
	 \loigiai{
		\begin{itemchoice}
			\itemch Không gian mẫu có $2^5 = 32$  phần tử.
			\itemch  Lần đầu xuất hiện mặt ngửa có $1$ lựa chọn, $4$ lần còn lại có $2^4 = 16$ cách.
			\itemch Biến cố đối $\overline{B}$ là \lq\lq Xuất hiện $5$ lần toàn mặt ngửa\rq\rq\,  có $1$ kết quả. Vậy $n(B) = 32 - 1 = 31$.
			\itemch Số lần mặt sấp hơn mặt ngửa nghĩa là số lần sấp  là $3$, $4$, hoặc $5$.\\
			Nên $n(C) =  \mathrm{C}_5^3 + \mathrm{C}_5^4 + \mathrm{C}_5^5 = 10 + 5 + 1 = 16$.
		\end{itemchoice}
	 }
\end{ex}

% nguồn: MaHoaNH_DungSai_Lop10
\begin{ex}%[0D0H1-2]%[ma:58281057]
	Gieo một đồng xu sau đó gieo một  con xúc xắc. Quan sát sự xuất hiện mặt sấp $(S)$, mặt ngửa $(N)$ của đồng xu và số chấm xuất hiện của con xúc xắc. Xét tính đúng sai của các khẳng định sau:
	\choiceTF 
		{\True Số phần tử không gian mẫu bằng $12$} 
		{Số phần tử của biến cố $A$: \lq\lq Đồng xu xuất hiện mặt sấp và con xúc xắc xuất  hiện mặt có số chấm chẵn\rq\rq\, bằng $2$}
		{Số phần tử của biến cố  $B$: \lq\lq Mặt ngửa của đồng xu và mặt có số chấm lẻ của con xúc xắc xuất hiện\rq\rq\, bằng $2$}
		{\True  Số phần tử của biến cố $C$: \lq\lq Mặt $6$ chấm xuất hiện\rq\rq\, bằng $2$}
	\loigiai{ 
		\begin{itemchoice}
			\itemch Không gian  mẫu là $\Omega = \{ S1, S2, S3, S4, S5, S6, N1, N2, N3, N4, N5, N6 \}$.\\
			Do đó, số phần  tử không gian mẫu là $n(\Omega) = 12$.
			\itemch Gọi $A$ là biến cố \lq\lq Đồng xu xuất hiện mặt sấp và con xúc xắc xuất  hiện mặt có số chấm chẵn\rq\rq\,.\\
			Ta có $A = \{ S2, S4, S6 \}$, do đó số phần tử của biến  cố này là $3$ chứ không phải $2$.
			\itemch Gọi $B$ là biến cố \lq\lq Mặt ngửa của đồng xu và mặt có số chấm lẻ của con xúc xắc xuất hiện\rq\rq\,.\\ 
			Ta có $B = \{ N1, N3, N5 \}$, do đó số phần tử  của biến cố này là $3$ chứ không phải $2$.
			 \itemch Gọi $C$ là biến cố \lq\lq Mặt $6$ chấm xuất hiện\rq\rq.\\
			Ta có $C =  \{ S6, N6 \}$, do đó số phần tử của biến cố này là $2$.
		\end{itemchoice}
	 }
\end{ex}

% nguồn: MaHoaNH_DungSai_Lop10
\begin{ex}%[0D0H1-2]%[ma:46e02225]
	Trên giá sách có  $4$ quyển sách toán, $3$ quyển sách lý, $2$ quyển sách hóa. Lấy ngẫu nhiên $3$ quyển sách. Xét tính đúng sai của các khẳng định sau:
	\choiceTF 
		 {\True $n(\Omega)=84$}
		{Số phần tử của biến cố $A$: \lq\lq Thuộc $3$ môn khác nhau\rq\rq\, bằng $20$} 
		{\True Số phần tử của biến cố  $B$: \lq\lq Đều là môn toán\rq\rq\, bằng $4$}
		{Số phần tử của biến cố $C$: \lq\lq Có ít nhất một quyển sách toán\rq\rq\, bằng $70$} 
	\loigiai{ 
		\begin{itemchoice}
			\itemch Không gian mẫu là kết quả của tổ hợp chập $3$ của $9$ phần tử, do đó số phần tử của không gian mẫu là $n(\Omega) =  \mathrm{C}_9^3 = 84$.
			\itemch Gọi  $A$ là biến cố \lq\lq Thuộc $3$ môn khác nhau\rq\rq.\\
			Ta có  $n(A) = 4 \cdot 3 \cdot 2 = 24$, chứ không phải $20$.
			\itemch Gọi $B$ là  biến cố \lq\lq Đều là môn toán\rq\rq.\\
			Ta có $n(B) = \mathrm{C}_4^3 = 4$, do đó số phần tử của biến cố này là $4$. 
			\itemch Gọi $C$ là biến cố \lq\lq Có ít nhất một quyển  sách toán\rq\rq.\\
			Gọi $\overline{C}$ là biến cố \lq\lq Không có quyển sách toán\rq\rq\,, tức là chọn $3$ quyển từ $5$ quyển sách còn lại (sách lý và hóa). Ta có $n(\overline{C}) = \mathrm{C}_5^3 =  10$.\\
			Do đó,  số phần tử của biến cố $C$ là $n(C) = n(\Omega) - n(\overline{C}) = 84 - 10 = 74$, chứ không phải $70$.
		\end{itemchoice}
	} 
\end{ex}

% nguồn: MaHoaNH_DungSai_Lop10
\begin{ex}%[0D0H1-2]%[ma:434fcba3]
	Xét phép thử là gieo một đồng xu gồm hai mặt sấp ngửa $3$ lần liên tiếp. Xét tính  đúng sai của các khẳng định sau:
	 \choiceTF
		{\True  $n(\Omega)=8$}
		 {\True Gọi $A$ là biến cố: \lq\lq Gieo được mặt sấp\rq\rq, khi đó $n\left(\overline{A}\right)=1$}
		{Gọi  $B$ là biến cố: \lq\lq Gieo được mặt sấp\rq\rq, khi đó $n(B)=1$}
		{\True Gọi $C$ là biến cố: \lq\lq Kết quả của lần gieo thứ hai  và thứ ba khác nhau\rq\rq, khi đó $n(C)=4$}
	 \loigiai{
		\begin{itemchoice}
			\itemch Ta có không gian mẫu  là $\Omega=\{SSS,SSN,SNS,NSS,SNN,NSN,NNS,NNN\}$.\\
			 Vậy $n(\Omega)=8$.
			\itemch Ta có $\overline{A}$: \lq\lq không gieo được mặt sấp\rq\rq. Suy ra  $\overline{A}=\{NNN\}$.\\
			 Vậy $n\left(\overline{A}\right)=1$.
			 \itemch Ta có $B=\{SSS,SSN,SNS,NSS,SNN,NSN,NNS\}$.\\
			 Vậy $n(B)=7$.
			\itemch Ta có  $C=\{SSN,SNS,NNS,NSN\}$.\\
			 Vậy $n(C)=4$.
		\end{itemchoice}
	 }
\end{ex}

% nguồn: MaHoaNH_DungSai_Lop10
\begin{ex}%[0D0H1-2]%[ma:5b165b7b]
	Một nhóm  có $6$ bạn nam và $5$ bạn nữ. Chọn ngẫu nhiên cùng một lúc ra $4$ bạn đi làm công tác tình nguyện. Xét tính đúng sai của các khẳng định sau:
	 \choiceTF
		{Số phần  tử của không gian mẫu là $320$}
		{\True Số các kết quả thuận lợi cho biến cố \lq\lq Trong $4$ bạn được chọn  có $2$ bạn nam và $2$ bạn nữ\rq\rq\, bằng $150$}
		{Số các kết quả thuận lợi cho biến cố \lq\lq Trong $4$  bạn được chọn có ít nhất $2$ bạn nữ\rq\rq\, bằng $225$}
		{Số các kết  quả thuận lợi cho biến cố \lq\lq Trong $4$ bạn được chọn có nhiều nhất $2$ bạn nữ\rq\rq\, bằng $260$}
	 \loigiai{
		\begin{itemchoice}
			\itemch Không gian mẫu có $\mathrm{C}_{11}^4 =  330$ phần tử.
			\itemch Ta có $\mathrm{C}_5^2 \cdot  \mathrm{C}_6^2 = 10 \cdot 15 = 150$.
			 \itemch Các trường hợp:
			\begin{itemize}
				\item $2$ nữ, $2$ nam: $\mathrm{C}_5^2 \cdot \mathrm{C}_6^2  = 150$.
				\item $3$ nữ, $1$ nam: $\mathrm{C}_5^3 \cdot \mathrm{C}_6^1 = 10 \cdot 6  = 60$.
				\item  $4$ nữ: $\mathrm{C}_5^4 = 5$.
			\end{itemize}
			Số các kết quả thuận lợi cho biến cố \lq\lq Trong $4$ bạn  được chọn có ít nhất $2$ bạn nữ\rq\rq\, bằng $150 + 60 + 5 = 215$.
			 \itemch Các trường hợp:
			\begin{itemize}
				\item $2$ nữ, $2$ nam: $150$. 
				\item $1$ nữ, $3$ nam: $\mathrm{C}_5^1 \cdot \mathrm{C}_6^3 = 5 \cdot 20 = 100$. 
				\item $0$ nữ,  $4$ nam: $\mathrm{C}_6^4 = 15$.
			\end{itemize}
			Số các kết quả thuận lợi cho biến cố \lq\lq Trong $4$ bạn được chọn có nhiều  nhất $2$ bạn nữ\rq\rq\, bằng $150 + 100 + 15 = 265$.
		\end{itemchoice}
	} 
\end{ex}


%%%%%%%%%%%%%%%% TRẢ LỜI NGẮN %%%%%%%%%%%%%%%%

% nguồn: MaHoaNH_TraLoiNgan_Lop10
\begin{bt}%[0D0N1-2]%[ma:b2d3dc69]
	Một nhóm bạn có $4$ bạn gồm $2$ bạn nam Mạnh, Dũng và hai nữ là Hoa, Lan được xếp ngẫu nhiên trên một  ghế dài. Kí hiệu MDHL là cách sắp xếp theo thứ tự: Mạnh, Dũng, Hoa, Lan. Tính số phần tử của không gian mẫu.
	\shortans{24} 
	\loigiai{ 
		Mỗi cách sắp xếp $4$ bạn vào $4$ chỗ ngồi là một hoán vị của  $4$ phần tử.\\
		Vì vậy số phần tử của không  gian mẫu là $4!=24$.
	 }
\end{bt}

% nguồn: MaHoaNH_TraLoiNgan_Lop10
\begin{bt}%[0D0H1-2]%[ma:ef322d8e]
	Có  $100$ tấm thẻ được đánh số từ $1$ đến $100$. Lấy ngẫu nhiên $5$ thẻ. Số phần tử của A: \lq\lq Số ghi trên các tấm thẻ được chọn là số chẵn\rq\rq\, có dạng $\mathrm{C}_m^5\,(m\in \mathbb{N};m\ge 5)$. Xác định giá trị của $m$.
	 \shortans{50}
	 \loigiai{
		Trong $100$ tấm thẻ thì có $50$  tấm thẻ là số chẵn và $50$ tấm thẻ là số lẻ.\\
		 Ta có $n(A)=\mathrm{C}_{50}^5$ nên khi đó $m=50$.
	} 
\end{bt}

% nguồn: MaHoaNH_TraLoiNgan_Lop10
\begin{bt}%[0D0H1-2]%[ma:1b272447]
	Hộp thứ nhất chứa $5$ quả bóng được đánh  số từ $1$ đến $5$. Hộp thứ hai chứa $6$ quả bóng được đánh số từ $1$ đến $6$. Chọn ngẫu nhiên mỗi hộp $1$ quả bóng. Có bao nhiêu kết quả thuận lợi cho biến cố \lq\lq Tổng các số ghi trên hai quả bóng không lớn hơn $8$\rq\rq?
	 \shortans{24}
	 \loigiai{
		Số kết quả có thể xảy ra của phép thử là $5 \times  6 = 30$.\\
		Vì số kết quả thuận  lợi cho biến cố \lq\lq Tổng các số ghi trên hai quả bóng lớn hơn $8$\rq\rq\,là $6$.\\
		Nên số kết  quả thuận lợi cho biến cố \lq\lq Tổng các số ghi trên hai quả bóng không lớn hơn $8$\rq\rq\,là $30 - 6 = 24$.
	} 
\end{bt}

% nguồn: MaHoaNH_TraLoiNgan_Lop10
\begin{bt}%[0D0H1-2]%[ma:ef43b2e7]
	Một hộp đựng $10$ viên bi xanh, $20$ viên bi đỏ, $15$ viên bi vàng. Chọn ngẫu nhiên cùng một lúc ra $3$  viên bi. Hãy xác định số các kết quả thuận lợi cho biến cố \lq\lq $3$ viên bi được chọn có màu khác nhau\rq\rq.
	\shortans{3,00} 
	\loigiai{ 
		Ta có  $\mathrm{C}_{10}^1$ cách chọn ra $1$ viên bi xanh, $\mathrm{C}_{20}^1$ cách chọn ra $1$ viên bi đỏ, $\mathrm{C}_{15}^1$ cách chọn ra $1$ viên bi vàng.\\
		Theo  quy tắc nhân ta có tất cả $\mathrm{C}_{10}^1\cdot \mathrm{C}_{20}^1\cdot \mathrm{C}_{15}^1=3\,000$ cách chọn ra $3$ viên bi có đủ $3$ màu.\\
		Do đó số kết quả thuận lợi cho biến cố \lq\lq $3$  viên bi được chọn có màu khác nhau\rq\rq\,là $3\,000$.
	} 
\end{bt}

% nguồn: MaHoaNH_TraLoiNgan_Lop10
\begin{bt}%[0D0H1-2]%[ma:6cb6b5ca]
	Hộp thứ nhất chứa $6$ quả bóng được đánh số từ $1$ đến $6$. Hộp thứ hai chứa $4$ quả bóng được đánh số từ  $1$ đến $4$. Chọn ngẫu nhiên mỗi hộp $1$ quả bóng. Có bao nhiêu kết quả thuận lợi cho biến cố \lq\lq Tổng các số ghi trên hai quả bóng không nhỏ hơn $5$\rq\rq?
	\shortans{18} 
	 \loigiai{
		Số kết quả có thể xảy ra của phép thử là $6 \times 4 =  24$.\\
		Vì số  kết quả thuận lợi cho biến cố \lq\lq Tổng các số ghi trên hai quả bóng nhỏ hơn $5$\rq\rq\, là $6$.\\
		Vậy số kết quả thuận lợi cho biến cố \lq\lq Tổng các số  ghi trên hai quả bóng không nhỏ hơn $5$\rq\rq\, là $24 - 6 = 18$.
	} 
\end{bt}

% nguồn: MaHoaNH_TraLoiNgan_Lop10
\begin{bt}%[0D0H1-2]%[ma:c0ba1124]
	Có $10$ học sinh nam và $10$ học sinh nữ xếp vào $1$ hàng dọc. Số các kết quả thuận lợi cho biến cố \lq\lq Học sinh nam và học  sinh nữ đứng xen kẽ nhau\rq\rq\, là $m$. Khi đó $\dfrac{m}{(9!)^2}$ bằng
	\shortans{200} 
	 \loigiai{
		Giả sử các vị trí đứng của $20$ bạn được đánh số từ $1$  đến $20$.
		\begin{itemize}
			\item Trường hợp $1$: Xếp $10$ bạn nam vào $10$ ô  chẵn, xếp $10$ bạn nữ vào $10$ ô lẻ. Ta có $10!10!$ cách.
			\item Trường hợp $2$: Xếp $10$ bạn nam vào $10$ ô lẻ, xếp $10$ bạn nữ vào  $10$ ô chẵn. Ta có $10!10!$ cách.
		\end{itemize}
		Theo  quy tắc cộng ta có số kết quả thuận lợi cho biến cố \lq\lq Học sinh nam và học sinh nữ đứng xen kẽ nhau\rq\rq\, là $m=10!10!+10!10!=2(10!)^2$.\\
		Do đó  $\dfrac{m}{(9!)^2}=200$.
	} 
\end{bt}

% nguồn: MaHoaNH_TraLoiNgan_Lop10
\begin{bt}%[0D0H1-2]%[ma:dc9e63f4]
	Một hộp đựng $15$ viên bi khác nhau gồm $4$ bi đỏ, $5$ bi trắng và $6$ bi vàng. Chọn ngẫu nhiên $4$ bi từ hộp, tính số phần tử của biến cố $X$: \lq\lq Chọn $4$ viên bi không  có đủ $3$ màu\rq\rq.
	 \shortans{645}
	\loigiai{ 
		Xét biến cố đối của $X$ là $\overline{X}$:  \lq\lq Chọn $4$ bi từ hộp có đủ ba màu\rq\rq:
		\begin{itemize}
			\item Trường hợp $1$: Chọn được $2$ bi đỏ, $1$ bi trắng, $1$  bi vàng, có $\mathrm{C}_4^2\cdot \mathrm{C}_5^1\cdot \mathrm{C}_6^1$ cách.
			\item Trường hợp $2$: Chọn được $1$ bi đỏ, $2$ bi trắng, $1$ bi vàng, có  $\mathrm{C}_4^1\cdot \mathrm{C}_5^2\cdot \mathrm{C}_6^1$ cách.
			\item Trường hợp $3$: Chọn được $1$ bi đỏ, $1$ bi trắng, $2$ bi vàng, có $\mathrm{C}_4^1\cdot \mathrm{C}_5^1\cdot \mathrm{C}_6^2$  cách.
		\end{itemize}
		Suy ra $n\left(\overline{X}\right)=\mathrm{C}_4^2\cdot \mathrm{C}_5^1\cdot \mathrm{C}_6^1+\mathrm{C}_4^1\cdot \mathrm{C}_5^2\cdot \mathrm{C}_6^1+\mathrm{C}_4^1\cdot  \mathrm{C}_5^1\cdot \mathrm{C}_6^2=720$.\\
		Số  phần tử của $X$ là $n(X)=n(\Omega)-n\left(\overline{X}\right)=\mathrm{C}_{15}^4-720=645$.
	} 
\end{bt}

% nguồn: MaHoaNH_TraLoiNgan_Lop10
\begin{bt}%[0D0V1-2]%[ma:bd762f59]
	Trong giải bóng đá nữ ở trường  THPT có $12$ đội tham gia, trong đó có hai đội của hai lớp $10A2$ và $10A5$. Ban tổ chức tiến hành bốc thăm ngẫu nhiên để chia thành hai bảng đấu $A$, $B$ mỗi bảng $6$ đội. Xác định số phần tử của biến cố để $2$ đội của hai lớp $10A2$ và $10A5$ ở cùng một bảng.
	 \shortans{420}
	 \loigiai{
		Gọi $A$ là biến cố: \lq\lq $2$ đội của hai lớp $10A2$ và $10A5$ ở cùng một bảng\rq\rq.\\ 
		Bốc $4$ đội  từ $10$ đội không tính hai lớp $10A2$ và $10A5$ vào bảng đã xếp hai đội của hai lớp $10A2$ và $10A5$; $6$ đội còn lại vào một bảng và hoán vị hai bảng.\\
		 Ta có $n(A)=\mathrm{C}_{10}^4\cdot 2!=420$.
	 }
\end{bt}

% nguồn: MaHoaNH_TraLoiNgan_Lop10
\begin{bt}%[0D0V1-2]%[ma:062e7ac8]
	Một nhóm có $4$ bạn nam và $6$ bạn nữ. Chọn ngẫu nhiên cùng một lúc ra $3$ bạn đi trực nhật. Hãy xác định số các  kết quả thuận lợi cho biến cố \lq\lq Trong $3$ bạn được chọn có ít nhất $1$ bạn nữ\rq\rq.
	\shortans{116} 
	 \loigiai{
		\begin{itemize}
			\item Trường hợp 1: Trong $3$ bạn được chọn có $1$ bạn nữ và $2$ bạn nam. Ta có $\mathrm{C}_6^1\cdot \mathrm{C}_4^2$ cách chọn. 
			\item Trường hợp 2: Trong $3$ bạn  được chọn có $2$ bạn nữ và $1$ bạn nam. Ta có $\mathrm{C}_6^2\cdot \mathrm{C}_4^1$ cách chọn.
			\item Trường hợp 3: Trong $3$ bạn được chọn có  $3$ bạn nữ. Ta có $\mathrm{C}_6^3$ cách chọn.
		\end{itemize}
		Theo quy  tắc cộng ta có số kết quả thuận lợi cho biến cố \lq\lq Trong $3$ bạn được chọn có ít nhất $1$ bạn nữ\rq\rq\, là $\mathrm{C}_6^1\cdot \mathrm{C}_4^2+\mathrm{C}_6^2\cdot \mathrm{C}_4^1+\mathrm{C}_6^3=116$.
	 }
\end{bt}

% nguồn: MaHoaNH_TraLoiNgan_Lop10
\begin{bt}%[0D0V1-2]%[ma:43aa3ff2]
	Một nhóm bạn có $4$ bạn gồm $2$ bạn nam Mạnh, Dũng và hai nữ là Hoa, Lan được xếp ngẫu nhiên trên một ghế dài. Kí hiệu MDHL là cách sắp xếp theo thứ tự: Mạnh, Dũng, Hoa, Lan. Tìm số phần tử của biến cố $B$: \lq\lq xếp nam và  nữ ngồi xen kẽ nhau\rq\rq.
	 \shortans{8}
	 \loigiai{
		\begin{itemize}
			\item Trường hợp $1$:  bạn nam ngồi đầu.\\
			Khi đó $2$ bạn nam xếp vào $2$ chỗ (số ghế $1$ và $3$)  có $2!$ cách, nữ xếp vào hai chỗ còn lại (ghế số $2$ và $4$) có $2!$ cách.\\
			Suy  ra số cách xếp là $2!\cdot 2!=4$ cách.
			 \item Trường hợp $2$: bạn nữ ngồi đầu. Tương tự có $4$ cách xếp.
		\end{itemize}
		Vậy theo quy tắc cộng số phần  tử của biến cố $N$ là $4+4=8$.
	 }
\end{bt}

% nguồn: MaHoaNH_TraLoiNgan_Lop10
\begin{bt}%[0D0V1-2]%[ma:0865d010]
	Cho tập hợp $A=\left\{1,2,3,4,5\right\}$. Gọi $S$ là tập hợp tất cả các số tự nhiên có ít nhất $3$ chữ số, các chữ số đôi một khác nhau được lập thành từ các chữ số thuộc  tập $A$. Chọn ngẫu nhiên một số từ $S$. Tính số phần tử của không gian mẫu trong phép thử trên.
	 \shortans{300}
	\loigiai{ 
		Ta tính số phần tử  thuộc tập $S$ như sau:
		\begin{itemize}
			 \item Số các số thuộc $S$ có $3$ chữ số là $\mathrm{A}_5^3$.
			\item Số các số thuộc $S$  có $4$ chữ số là $\mathrm{A}_5^4$.
			\item Số các  số thuộc $S$ có $5$ chữ số là $\mathrm{A}_5^5$.
		\end{itemize}
		Suy ra số  phần tử của tập $S$ là $\mathrm{A}_5^3+\mathrm{A}_5^4+\mathrm{A}_5^5=300$.\\
		Phép thử \lq\lq Chọn ngẫu nhiên $1$ số từ tập  $S$\rq\rq. \\
		Suy ra số phần tử của không gian mẫu  là $n(\Omega)=\mathrm{C}_{300}^1=300$.
	} 
\end{bt}

% nguồn: MaHoaNH_TraLoiNgan_Lop10
\begin{bt}%[0D0V1-2]%[ma:3ce98ff2]
	Một hộp đựng $5$ viên bi xanh, $4$ viên bi đỏ, $7$  viên bi vàng. Chọn ngẫu nhiên cùng một lúc ra $8$ viên bi. Số các kết quả thuận lợi cho biến cố \lq\lq $8$ viên bi được chọn có đủ $3$ màu\rq\rq\, là $n$. Khi đó $\dfrac{n}{3}$ bằng
	\shortans{4,07} 
	 \loigiai{
		\begin{itemize}
			\item Trường hợp 1: Trong $8$ viên bi được chọn chỉ có $2$ màu xanh và  đỏ. Ta có $\mathrm{C}_9^8$ cách chọn.
			\item Trường hợp 2: Trong $8$ viên bi được chọn chỉ có $2$ màu xanh và vàng. Ta có $\mathrm{C}_{12}^8$ cách chọn. 
			\item Trường hợp 3: Trong $8$ viên bi được chọn chỉ có $2$ màu đỏ và vàng. Ta có $\mathrm{C}_{11}^8$ cách  chọn.\\
			Chọn ra $8$ viên bi bất kỳ ta có $\mathrm{C}_{16}^8 = 12\,870$  cách chọn.
		\end{itemize}
		Chọn ra $8$ viên bi không có đủ $3$ màu ta có $\mathrm{C}_9^8 + \mathrm{C}_{12}^8 + \mathrm{C}_{11}^8 = 669$ cách chọn.\\ 
		Chọn ra $8$ viên bi có đủ $3$ màu ta có $12\,870 - 669 =  12\,201$ cách chọn.\\
		Do đó số kết quả thuận lợi cho biến cố \lq\lq $8$ viên bi được  chọn có đủ $3$ màu\rq\rq\, là $n=12\,201$.\\
		Vậy $\dfrac{n}{3}=4\,067$. 
	} 
\end{bt}


%%%%%%%%%%%%%%%% TỰ LUẬN %%%%%%%%%%%%%%%%

% nguồn: MaHoaNH_TuLuan_Lop10
\begin{bt}%[0D0H1-2]%[ma:74395f7a]
	Khi ta tung  một đồng xu có 2 mặt, ta hoàn toàn không biết trước được kết quả của nó, tuy nhiên ta lại biết chắc chắn rằng đồng xu rơi xuống sẽ ở một trong 2 trạng thái: sấp $(S)$ hoặc ngửa $(N)$.\\
	Khi đó không  gian mẫu của phép thử là $\Omega=\{S; N\}$.
\end{bt}

% nguồn: MaHoaNH_TuLuan_Lop10
\begin{bt}%[0D0H1-2]%[ma:3d28981d]
	Chọn ngẫu nhiên một gia đình có ba con ($12$ tuổi, $15$  tuổi và $18$ tuổi). Quan sát giới tính của ba người con này. Mô tả không gian mẫu.
	\loigiai{ 
		Chọn ngẫu nhiên một gia đình có ba con và quan sát giới tính của ba người con này, ta có không gian mẫu $\Omega = \{TTT; TTG; TGT; TGG; GTT; GTG; GGT; GGG\}$ nên số phần tử của không gian mẫu là $n(\Omega) =  8$.
	} 
\end{bt}

% nguồn: MaHoaNH_TuLuan_Lop10
\begin{bt}%[0D0H1-2]%[ma:95c1256f]
	Chọn ngẫu nhiên hai số khác nhau từ $5$  số nguyên dương đầu tiên. Mô tả không gian mẫu.
	\loigiai{ 
		Không gian  mẫu $\Omega = \{(1; 2); (1; 3); (1; 4); (1; 5); (2; 3); (2; 4); (2; 5); (3; 4); (3; 5); (4; 5)\}$.
	} 
\end{bt}

% nguồn: MaHoaNH_TuLuan_Lop10
\begin{bt}%[0D0H1-2]%[ma:6156391f]
	Xét phép thử:  \lq\lq Tung một đồng xu hai lần liên tiếp\rq\rq. Sự kiện: \lq\lq Kết quả hai lần gieo khác nhau\rq\rq\ tương ứng với biến cố nào của phép thử trên.\\
	 \loigiai{
		Sự kiện: \lq\lq Kết quả  hai lần gieo khác nhau\rq\rq\ tương ứng với biến cố $A=\{S N; N S\}$.
	} 
\end{bt}

% nguồn: MaHoaNH_TuLuan_Lop10
\begin{bt}%[0D0H1-2]%[ma:80108c2f]
	Gieo  một con xúc xắc hai lần liên tiếp. Hãy mô tả không gian mẫu của phép thử.\\
	 \loigiai{
		Không gian mẫu của phép thử  $\Omega=\{(i; j) \mid i=\overline{1,6}, j=\overline{1,6}\}$.
	 }
\end{bt}

% nguồn: MaHoaNH_TuLuan_Lop10
\begin{bt}%[0D0H1-2]%[ma:61416175]
	Chọn ngẫu nhiên một số nguyên dương không lớn  hơn $30$. Mô tả không gian mẫu.
	\loigiai{ 
		Không gian mẫu $\Omega =  \{1; 2; 3; \ldots; 27; 28; 29; 30\}$.
	} 
\end{bt}

% nguồn: MaHoaNH_TuLuan_Lop10
\begin{bt}%[0D0H1-2]%[ma:cdd93928]
	Gieo đồng thời một con xúc xắc và một  đồng xu. Mô tả không gian mẫu.
	 \loigiai{
		Kí  hiệu $S$ và $N$ tương ứng là đồng xu ra mặt sấp và đồng xu ra mặt ngửa. Khi đó không gian mẫu là
		$$\Omega = \{(1; S); (2; S); (3; S); (4; S); (5; S); (6; S); (1; N); (2; N); (3; N); (4;  N); (5; N); (6; N)\}.$$
	 }
\end{bt}

% nguồn: MaHoaNH_TuLuan_Lop10
\begin{bt}%[0D0H1-2]%[ma:ba58ed58]
	Gieo đồng thời ba đồng xu cân đối đồng  chất. Mô tả không gian mẫu.
	 \loigiai{
		Kí hiệu $S$ và $N$ tương ứng là đồng xu ra mặt sấp  và đồng xu ra mặt ngửa. Khi đó không gian mẫu là
		\[\Omega = \{(S; S; S); (S; S; N); (S; N; S); (S; N; N); (N; S; S); (N; S; N); (N; N; S); (N; N; N)\}.\] 
	 }
\end{bt}

% nguồn: MaHoaNH_TuLuan_Lop10
\begin{bt}%[0D0H1-2]%[ma:4e032d81]
	Nhóm $1$ trong văn phòng có $1$ nhân viên nữ là Xuân, $3$ nhân viên nam là Hạ, Thu, Đông. Quản lý chọn ngẫu nhiên $2$ nhân viên $1$ nam, $1$ nữ để phỏng vấn. Phép thử ngẫu nhiên là  gì? Mô tả không gian mẫu.
	 \loigiai{
		Phép thử  ngẫu nhiên là chọn $2$ nhân viên $1$ nam, $1$ nữ để phỏng vấn.\\
		Không gian mẫu là  $\Omega = \{(\text{Xuân, Hạ}); (\text{Xuân, Thu}); (\text{Xuân, Đông})\}$.
	 }
\end{bt}

% nguồn: MaHoaNH_TuLuan_Lop10
\begin{bt}%[0D0H1-2]%[ma:9a3ba8db]
	Xếp ba người ngồi thành hàng ngang. Mô tả không gian mẫu  của phép thử đó.\\
	 \loigiai{
		Đặt tên  ba người theo thứ tự là $A$, $B$, $C$. \\
		Khi đó,  không gian mẫu của phép thử là
		$\Omega=\{A B C; A C B; B A C; B C A; C A  B; C B A\}$.
	 }
\end{bt}

% nguồn: MaHoaNH_TuLuan_Lop10
\begin{bt}%[0D0H1-2]%[ma:8163de0c]
	Đề thi môn Toán có $50$ câu trắc nghiệm, mỗi câu có  bốn phương án $A;\ B;\ C;\ D$ trong đó có một phương án đúng. Bạn An làm bài bằng cách chọn ngẫu nhiên một phương án trong bốn phương án của mỗi câu. Hãy xác định số phần tử của không gian mẫu.
	 \loigiai{
		 Phép thử: \lq\lq Chọn ngẫu nhiên cả $50$ câu trả lời, mỗi câu một phương án \rq\rq.\\
		 Mỗi câu có $4$ phương án trả lời nên có $4^{50}$ khả năng.\\
		Vậy số phần  tử của không gian mẫu là $n(\Omega)=4^{50}$.
	 }
\end{bt}

% nguồn: MaHoaNH_TuLuan_Lop10
\begin{bt}%[0D0H1-2]%[ma:bd8f5aba]
	Tung một đồng tiền năm lần liên tiếp. Tính số phần tử của biến  cố $D$: \lq\lq Mặt sấp xuất hiện ít nhất một lần\rq\rq.\\
	 \loigiai{
		Số phần tử của không gian mẫu  là $n_\Omega=2^5=32$.\\
		$\overline{D}$: \lq\lq Cả năm lần gieo đều xuất hiện mặt  ngửa\rq\rq. Khi đó $\overline{D}=\{N N N N N\}$ nên $n(\overline{D})=1$.\\
		Do đó $n(D)=32 - 1 =  31$.
	 }
\end{bt}

% nguồn: MaHoaNH_TuLuan_Lop10
\begin{bt}%[0D0H1-2]%[ma:5b5efd16]
	Xét phép thử: \lq\lq Gieo một con xúc xắc hai lần liên tiếp\rq\rq. Sự kiện \lq\lq Số chấm xuất hiện  ở lần gieo thứ nhất là $2$ \rq\rq\ tương ứng với biến cố nào của phép thử trên.\\
	 \loigiai{
		Sự kiện \lq\lq Số chấm xuất hiện ở lần gieo thứ nhất là $2$ \rq\rq\ tương  ứng với biến cố
		 \[
		B=\{(2; 1);(2;  2);(2; 3);(2; 4);(2; 5);(2; 6)\}.
		 \]
	 }
\end{bt}

% nguồn: MaHoaNH_TuLuan_Lop10
\begin{bt}%[0D0H1-2]%[ma:7cdbe15d]
	Trong một hộp đựng bi có $5$ viên bi được  đánh số $1;\ 2;\ 3;\ 4;\ 5$. Chọn ngẫu nhiên hai viên bi trong hộp, hỏi trong phép thử này số phần tử của không gian mẫu là bao nhiêu? Tìm không gian mẫu.
	 \loigiai{
		Phép  thử: \lq\lq Chọn ngẫu nhiên $2$ viên bi từ $5$ viên bi \rq\rq.\\
		 Số phần tử của không gian mẫu là $n(\Omega)=\mathrm{C}_5^2=10$.\\
		Không gian mẫu: 
		 \[
		\Omega=\{\{1;\ 2\};\ \{1;\ 3\};\  \{1;\ 4\};\ \{1;\ 5\};\ \{2;\ 3\};\ \{2;\ 4\};\ \{2;\ 5\};\ \{3;\ 4\};\ \{3;\ 5\};\ \{4;\ 5\}\}.
		\] 
	 }
\end{bt}

% nguồn: MaHoaNH_TuLuan_Lop10
\begin{bt}%[0D0H1-2]%[ma:8dd440e8]
	 Một quán ăn vặt có các món chè là chè bưởi, chè đậu xanh và chè thập cẩm, các món kem là kem xôi và kem sôcôla. Một thực khách vào quán và chọn ngẫu nhiên một món trong các món trên, hỏi trong phép thử này số phần tử của không gian mẫu là bao nhiêu? Tìm không gian mẫu.
	 \loigiai{
		Phép thử: \lq\lq Chọn một món ăn có sẵn \rq\rq. 
		\begin{itemize}
			\item  Chọn món chè, có $3$ cách.
			 \item Chọn món kem, có $2$ cách.
		\end{itemize}
		Theo quy tắc cộng thì  số phần tử của không gian mẫu là $n(\Omega)=5$.\\
		Không gian mẫu  $\Omega=\{\text{chè bưởi};\ \text{chè đậu xanh};\ \text{chè thập cẩm};\ \text{kem xôi};\ \text{kem sô-cô-la}\}$.
	} 
\end{bt}

% nguồn: MaHoaNH_TuLuan_Lop10
\begin{bt}%[0D0H1-2]%[ma:5d7c6335]
	Xét phép thử gieo một con xúc xắc ba lần liên tiếp. 
	 \[
	 C=\{(1; 3; 5);(1; 5; 3);(3; 1; 5);(3; 5; 1);(5; 1; 3);(5; 3; 1)\}.
	\] 
	Phát biểu biến cố của không gian mẫu (trong phép thử trên) dưới dạng  mệnh đề nêu sự kiện.
	 \\
	\loigiai{ 
		Tập con $C$ bao gồm tất cả các phần tử của không gian mẫu có tính chất đặc trưng là số chấm xuất hiện của ba lần gieo khác  nhau và đều là số lẻ.\\
		Vậy biến cố $C$ có thể phát biểu dưới dạng mệnh đề nêu sự kiện: \lq\lq Số chấm  xuất hiện trong cả ba lần gieo khác nhau và đều là số lẻ\rq\rq.
	 }
\end{bt}

% nguồn: MaHoaNH_TuLuan_Lop10
\begin{bt}%[0D0H1-2]%[ma:00548e18]
	Phần thưởng ở lớp trong dịp thi đua điểm tốt là bút mực, bút bi, bút chì, vở, thước kẻ, compa. Bạn Hoa đạt nhiều điểm tốt  nên được tham gia chọn $1$ phần quà.
	\begin{enumerate}
		\item Mô tả  không gian mẫu.
		 \item Gọi $A$ là biến cố \lq\lq Bạn Hoa chọn được một món quà có thể viết\rq\rq. $A$ là tập con nào của không gian mẫu?
	\end{enumerate}
	 \loigiai{
		\begin{enumerate}
			\item Không gian mẫu là $\Omega=\{$ bút mực, bút bi, bút chì, vở, thước kẻ,  compa $\}$.
			\item $A=\{$ bút mực, bút bi,  bút chì $\}$.
		\end{enumerate}
	} 
\end{bt}

% nguồn: MaHoaNH_TuLuan_Lop10
\begin{bt}%[0D0H1-2]%[ma:921f90b4]
	Có $12$ tấm thẻ được đánh số từ $1$ đến $12$. Rút ngẫu nhiên đồng thời $5$ tấm thẻ. Tính số phần tử của biến cố: \lq\lq Tổng các số ghi trên $5$ tấm thẻ được rút ra  là một số lẻ\rq\rq.
	\loigiai{ 
		 Gọi $A$ là biến cố \lq\lq năm thẻ lấy ra có tổng các số ghi trên năm thẻ là số lẻ\rq\rq.\\
		Ta có, trong $12$ tấm thẻ được  đánh số từ $1$ đến $12$ thì có $6$ tấm thẻ ghi số chẵn và $6$ tấm thẻ ghi số lẻ.\\
		Để tổng các số ghi trên $5$ tấm  thẻ được lấy ra là số lẻ thì số thẻ ghi số lẻ phải là lẻ.\\
		Khi đó  ta có
		\begin{itemize}
			\item Chọn $1$ thẻ lẻ và $4$ thẻ chẵn: $\mathrm{C}_6^1 \cdot  \mathrm{C}_6^4 = 90$.
			\item Chọn $3$ thẻ lẻ và $2$ thẻ chẵn: $\mathrm{C}_6^3 \cdot  \mathrm{C}_6^2 = 300$.
			\item Chọn  $5$ thẻ đều lẻ: $\mathrm{C}_6^5 = 6$.
		\end{itemize}
		 Vậy $n(A) = 90 + 300 + 6 = 396$.
	} 
\end{bt}

% nguồn: MaHoaNH_TuLuan_Lop10
\begin{bt}%[0D0H1-2]%[ma:83067db5]
	Có hai hộp đựng bi, hộp thứ nhất đựng ba viên bi có màu lần lượt là lam, đỏ, vàng; hộp thứ hai có đựng bốn viên bi có màu lần lượt là tím, trắng, lục, cam. Bạn Khoa lấy ngẫu nhiên mỗi hộp một viên bi, hỏi trong phép thử  này số phần tử của không gian mẫu là bao nhiêu? Tìm không gian mẫu.
	\loigiai{ 
		Phép thử \lq\lq  Lấy ngẫu nhiên mỗi hộp một viên bi \rq\rq.
		\begin{itemize}
			\item Lấy $1$ bi từ hộp thứ nhất có  $3$ cách chọn.
			\item Lấy $1$ bi từ hộp  thứ hai có $4$ cách chọn.
		\end{itemize}
		 Theo quy tắc nhân, số phần tử của không gian mẫu là $n(\Omega)=3 \cdot 4 = 12$.\\
		Không gian mẫu 
		\begin{eqnarray*}
			\Omega=\{(\text{lam};\ \text{tím});\ (\text{lam};\ \text{trắng});\ (\text{lam};\ \text{lục});\ (\text{lam};\ \text{cam});\  (\text{đỏ};\ \text{tím});\ (\text{đỏ};\ \text{trắng});\ (\text{đỏ};\ \text{lục});\ (\text{đỏ};\ \text{cam});\\ (\text{vàng};\ \text{tím});\ (\text{vàng};\ \text{trắng});\ (\text{vàng};\ \text{lục});\ (\text{vàng};\ \text{cam})\}.
		\end{eqnarray*}
	} 
\end{bt}

% nguồn: MaHoaNH_TuLuan_Lop10
\begin{bt}%[0D0H1-2]%[ma:45e48d5a]
	Có $3$ bông hoa: hồng, vàng, trắng trên bàn. Người ta lấy  ngẫu nhiên $2$ bông để cắm vào $2$ lọ: một lọ pha lê và một lọ gốm. Hỏi trong phép thử này số phần tử của không gian mẫu là bao nhiêu? Tìm không gian mẫu.
	\loigiai{ 
		Phép thử: \lq\lq Chọn ngẫu nhiên $2$ bông  hoa từ $3$ bông rồi cắm vào $2$ lọ \rq\rq.\\
		Số phần tử của không gian  mẫu là $n(\Omega)=\mathrm{A}_3^2=6$.\\
		 Không gian mẫu:
		\begin{eqnarray*}
			 \Omega=\{
			 (\text{hoa hồng - lọ pha lê};\ \text{hoa vàng - lọ gốm});\
			(\text{hoa vàng  - lọ pha lê};\ \text{hoa hồng - lọ gốm}); \\
			(\text{hoa  hồng - lọ pha lê}; \text{hoa trắng - lọ gốm});
			(\text{hoa trắng - lọ pha lê};\ \text{hoa hồng -  lọ gốm}); \\
			(\text{hoa vàng - lọ pha lê};\ \text{hoa trắng - lọ gốm});\ 
			(\text{hoa trắng - lọ pha lê};\ \text{hoa vàng  - lọ gốm})
			 \}.
	\end{eqnarray*}}
\end{bt}

% nguồn: MaHoaNH_TuLuan_Lop10
\begin{bt}%[0D0H1-2]%[ma:8c358ca6]
	Bạn Lan có $10$ chiếc thẻ, mỗi chiếc được đánh một số tự nhiên từ $1$ đến $10$. Bạn Lan chọn ngẫu nhiên ra hai chiếc thẻ trong  số đó.
	\begin{enumerate}
		\item Hãy  xác định số phần tử của không gian mẫu.
		\item Hãy xác định  các phần tử của biến cố \lq\lq Tổng các số trên hai chiếc thẻ bằng $5$\rq\rq.
		\item Hãy xác định số phần tử của biến cố \lq\lq Tổng các  số trên hai chiếc thẻ là số lẻ\rq\rq.
		\item Hãy xác định số  phần tử của biến cố \lq\lq Tích các số trên hai tấm thẻ là số lẻ\rq\rq.
	\end{enumerate}
	\loigiai{ 
		\begin{enumerate}
			\item Số phần tử của không  gian mẫu là $n(\Omega) = \mathrm{C}_{10}^2 = 45$.
			\item Gọi $A$ là biến cố \lq\lq Tổng các số trên hai chiếc thẻ bằng $5$\rq\rq.\\ 
			Ta có $A = \{\{1; 4\}, \{2; 3\}\} \Rightarrow n(A) = 2$. 
			\item Tổng của hai số là số lẻ khi một số  chẵn và một số lẻ được chọn.\\
			Có $5$ số  lẻ: $1$, $3$, $5$, $7$, $9$ và $5$ số chẵn: $2$, $4$, $6$, $8$, $10$.\\
			 Suy ra số phần tử của biến cố là $\mathrm{C}_5^1 \cdot \mathrm{C}_5^1 = 25$.
			\item Tích là số lẻ  khi cả hai số đều là số lẻ.\\
			Số phần tử của  biến cố là $\mathrm{C}_5^2 = 10$.
		\end{enumerate}
	} 
\end{bt}

% nguồn: MaHoaNH_TuLuan_Lop10
\begin{bt}%[0D0H1-2]%[ma:234c80f7]
	Thực hiện phép thử  gieo một đồng xu cân đối và đồng chất $3$ lần liên tiếp.
	\begin{enumerate}
		\item Hãy mô tả không gian mẫu và xác định  số phần tử của không gian mẫu.
		\item Hãy xác định các  phần tử của biến cố \lq\lq Mặt sấp xuất hiện đúng một lần\rq\rq.
		 \item Hãy xác định các phần tử của biến cố \lq\lq Mặt sấp xuất hiện lần thứ hai\rq\rq.
		 \item Hãy xác định các phần tử của biến cố \lq\lq Mặt ngửa xuất hiện ít nhất một lần\rq\rq.
	\end{enumerate}
	\loigiai{ 
		\begin{enumerate}
			\item Kí hiệu mặt sấp là $S$, mặt ngửa là $N$.\\ 
			Ta có  không gian mẫu
			\[\Omega = \{SSS;\ SSN;\ SNS;\  NSS;\ SNN;\ NSN;\ NNS;\ NNN\} \Rightarrow n(\Omega) = 8.\]
			\item Gọi $A$ là biến cố \lq\lq Mặt sấp xuất hiện đúng một lần\rq\rq.\\ 
			Khi đó $A = \{SNN;\  NSN;\ NNS\} \Rightarrow n(A) = 3$.
			\item Gọi $B$ là  biến cố \lq\lq Mặt sấp xuất hiện lần thứ hai\rq\rq.\\
			Ta có $B = \{NSN;\ SSN;\  NSS;\ SSS\} \Rightarrow n(B) = 4$.
			\item Gọi $C$ là biến  cố \lq\lq Mặt ngửa xuất hiện ít nhất một lần\rq\rq.\\
			Khi đó $C = \{SSN;\ SNS;\ NSS;\ SNN;\ NSN;\ NNS;\  NNN\} \Rightarrow n(C) = 7$.
		\end{enumerate}
	} 
\end{bt}

% nguồn: MaHoaNH_TuLuan_Lop10
\begin{bt}%[0D0H1-2]%[ma:400cd86c]
	Trong một hộp có $4$ quả  cầu màu xanh, $3$ quả cầu màu đỏ và $5$ quả cầu màu vàng (các quả cầu đôi một khác nhau). Chọn ngẫu nhiên ra $3$ quả cầu.
	\begin{enumerate}
		\item Hãy xác định số phần tử  của không gian mẫu.
		\item Hãy xác định số phần  tử của biến cố: \lq\lq Ba quả cầu được lấy ra có đủ cả ba màu\rq\rq.
		\item Hãy xác định số phần tử của biến cố: \lq\lq Có đúng một quả cầu màu đỏ được  lấy ra\rq\rq.
	\end{enumerate}
	 \loigiai{
		\begin{enumerate}
			\item Tổng số quả  cầu là $4 + 3 + 5 = 12$.\\
			Số phần tử của không gian mẫu là $n(\Omega) = \mathrm{C}_{12}^3  = 220$.
			\item Để có đủ ba màu, ta chọn $1$ quả xanh, $1$ quả đỏ, $1$ quả vàng 
			\begin{itemize}
				\item Chọn $1$ quả xanh có $4$  cách.
				\item Chọn $1$ quả  đỏ có $3$ cách.
				 \item Chọn $1$ quả đỏ có $5$ cách.
			\end{itemize}
			Theo quy tắc nhân, ta có $4 \cdot 3 \cdot 5 =  60$ cách.
			\item Có đúng $1$ quả đỏ, nên chọn $1$ trong $3$ quả đỏ và $2$ trong $9$ quả còn lại (xanh  + vàng)
			\begin{itemize}
				\item Chọn $1$  quả đỏ có $3$ cách.
				\item Chọn $2$ trong $9$ quả  xanh hoặc vàng có $\mathrm{C}_9^2=36$ cách.
			\end{itemize}
			Theo quy tắc nhân,  ta có $3 \cdot 36=108$ cách.
		\end{enumerate}
	} 
\end{bt}

% nguồn: MaHoaNH_TuLuan_Lop10
\begin{bt}%[0D0H1-2]%[ma:5b549d43]
	Lớp $10A$ có $45$ học sinh, trong đó có $20$ học sinh nam và $25$ học sinh nữ. Cô  giáo chủ nhiệm lớp chọn ngẫu nhiên $2$ bạn vào ban cán sự lớp.
	\begin{enumerate}
		\item Hãy xác định số phần tử của không  gian mẫu.
		\item Hãy xác định số phần tử của biến  cố \lq\lq Có một học sinh nam và một học sinh nữ được chọn vào ban cán sự lớp\rq\rq.
		\item Hãy xác định số phần tử của biến cố \lq\lq Hai học  sinh được chọn là những học sinh nữ\rq\rq.
		\item Hãy xác định số phần tử của biến cố \lq\lq Hai học  sinh được chọn là những học sinh nam\rq\rq.
	\end{enumerate}
	\loigiai{ 
		\begin{enumerate}
			\item Số  phần tử của không gian mẫu là $n(\Omega) = \mathrm{C}_{45}^2 = 990$.
			\item Số cách chọn một học sinh nam  là $\mathrm{C}_{20}^1$, một học sinh nữ là $\mathrm{C}_{25}^1$.\\
			 Do đó, số phần tử của biến cố cần tìm là
			 \[
			\mathrm{C}_{20}^1 \cdot \mathrm{C}_{25}^1 =  20 \cdot 25 = 500.
			 \]
			\item Số phần tử của biến cố \lq\lq Hai học sinh được chọn là những học sinh nữ\rq\rq\  là
			\[ 
			 \mathrm{C}_{25}^2 = \dfrac{25 \cdot 24}{2} = 300.
			\] 
			\item Số phần tử của biến  cố \lq\lq Hai học sinh được chọn là những học sinh nam\rq\rq\ là
			 \[
			\mathrm{C}_{20}^2 = \dfrac{20 \cdot 19}{2}  = 190.
			 \]
		\end{enumerate}
	 }
\end{bt}

% nguồn: MaHoaNH_TuLuan_Lop10
\begin{bt}%[0D0H1-2]%[ma:74a17f3a]
	 Gieo một con xúc xắc cân đối và đồng chất hai lần liên tiếp.
	\begin{enumerate}
		\item  Hãy mô tả không gian mẫu và tính số phần tử của không gian mẫu của phép thử này.
		\item Hãy xác định các phần tử của biến cố \lq\lq Mặt $5$ chấm xuất hiện lần gieo đầu tiên\rq\rq. 
		\item Hãy xác định các phần  tử của biến cố \lq\lq Tổng số chấm xuất hiện sau $2$ lần gieo bằng $8$\rq\rq.
		\item Hãy xác định các phần tử của biến cố \lq\lq Số chấm xuất hiện ở hai lần gieo là như nhau\rq\rq. 
	\end{enumerate}
	 \loigiai{
		\begin{enumerate}
			\item Không gian mẫu là $\Omega = \{(i;  j) \mid i,j = \overline{1,6}\}$.\\
			 Số phần tử của không gian mẫu là $n(\Omega) = 6 \cdot 6 = 36$.
			\item Gọi $A$  là biến cố \lq\lq Mặt $5$ chấm xuất hiện lần gieo đầu tiên\rq\rq.\\
			 Ta có
			 \[
			A = \{(5; 1);\ (5; 2);\ (5; 3);\  (5; 4);\ (5; 5);\ (5; 6)\} \Rightarrow n(A) = 6.
			 \]
			\item Gọi $B$  là biến cố \lq\lq Tổng số chấm xuất hiện sau hai lần gieo bằng $8$\rq\rq.\\
			Ta có 
			\[ 
			B = \{(2; 6);\ (3; 5);\ (4; 4);\ (5; 3);\ (6; 2)\} \Rightarrow  n(B) = 5.
			\] 
			\item Gọi $C$ là  biến cố \lq\lq Số chấm xuất hiện ở hai lần gieo là như nhau\rq\rq.\\
			Khi đó: 
			 \[
			C = \{(1; 1);\ (2; 2);\ (3; 3);\  (4; 4);\ (5; 5);\ (6; 6)\} \Rightarrow n(C) = 6.
			 \]
		\end{enumerate}
	 }
\end{bt}
